\section{Энергия: машина платит за обучение, не за счёт}
\label{sec:energy}

\subsection{Где тратит фон-неймановская машина}

После конца деннардовского масштабирования энергия вычисления доминируется
\emph{перемещением} операндов, не арифметикой: операция с плавающей точкой стоит
$\sim\!1$~пДж, внутрикристальное слово $\sim\!1$~пДж, а внекристальное DRAM-слово
$\sim\!20$~пДж --- на один-два порядка больше и растёт с расстоянием.
Обучающаяся машина усугубляет это: тренировка гоняет активации и градиенты между
счётом и памятью на каждом шаге. Фон-неймановский энергетический пол ---
\emph{перемещение данных}, и он лежит на семь-восемь порядков выше
термодинамического (Ландауэрова) предела~\cite{landauer1961,bennett1973} того же
вычисления.

\subsection{Где тратит \TALOS{}: новый учёт}

Единственная инструкция \TALOS{} чисто расщепляется на обратимую часть и две
экспортирующие энтропию части (Таблица~\ref{tab:three-terms}), и это расщепление
индуцирует энергетический закон без аналога в слоистой машине.

\begin{claimbox}[Закон энергия--обучение \status{Т} для обратимого ядра, \status{С} для коэффициента]
Унитарный член $-i[H_{\mathrm{eff}},\Gam]$ --- \emph{изометрия} всякой монотонной
метрики (T-262): он сохраняет фон-неймановскую энтропию состояния, поэтому по
Ландауэру стоит \emph{ноль} свободной энергии. Весь расход свободной энергии
относится на два экспортирующих энтропию члена:
\begin{equation}
\label{eq:energy-law}
E_{\mathrm{tick}} \;\geq\; kT\ln 2 \,\cdot\, \Delta H_{\mathrm{exported}},
\qquad
\Delta H_{\mathrm{exported}} \;=\; \underbrace{\Delta H_{\Regen}}_{\text{обучение}}
\;+\; \underbrace{\Delta H_{\Diss}^{(+)}}_{\text{поддержание}},
\end{equation}
где $\Delta H_{\Regen}$ --- порядок, который создаёт регенерация
(дорогое-по-свободной-энергии направление, $V_{\mathrm{hed}}=dP/d\tau>0$), а
$\Delta H_{\Diss}^{(+)}$ --- лишь та энтропия, которую надо ре-экспортировать,
чтобы удержать состояние против утечки диссипатора. \emph{Счёт бесплатен; лишь
забывание-в-порядок (поддержание) и обучение стоят энергии.}
\end{claimbox}

Два предельных случая делают закон конкретным.
\begin{itemize}[leftmargin=1.5em,itemsep=1pt]
  \item \textbf{Холон в своей неподвижной точке} ($\Gam=\rhostar$, сошёлся, не
        учится) имеет $\Delta H_{\Regen}=0$; его единственная цена --- член
        поддержания, свободная энергия для удержания $\rhostar$ против релаксации.
        Если задача не требует забывания, диссипатор можно погасить, и машина
        считает \emph{обратимо на полу Ландауэра} --- чистое стирание ноль.
  \item \textbf{Холон, который учится,} тратит пропорционально тому, как быстро
        он движется к своей самомодели, $\propto (dP/d\tau)_+$. Мощность
        отслеживает скорость обучения прямо: машина, которая научилась, потребляет
        меньше машины, которая учится, без изменения того, что она считает.
\end{itemize}

\subsection{Числа}

При $T=300$~К $kT\ln 2 = 2{,}87\times10^{-21}$~Дж на необратимо экспортированный
бит ($0{,}018$~эВ). Тик, ограниченный обучением, перекраивающий несколько бит
$34$-параметрического состояния, стоит порядка $10^{-20}$~Дж; худший случай ---
ре-экспорт всей энтропии состояния каждый тик, чего вдали от холодного старта не
бывает --- есть $48\,kT\ln 2 = 1{,}4\times10^{-19}$~Дж. Против одного
внекристального доступа к памяти фон-неймановской машины
($\sim\!2\times10^{-11}$~Дж) ограниченный обучением тик \TALOS{} сидит примерно
на \textbf{девять порядков} ниже. Архитектура не \emph{достигает} пола Ландауэра
в реальном субстрате --- накладные интерфейса, управления и утечки сидят выше
него --- но она \emph{убирает член перемещения данных, задающий фон-неймановский
пол}, оставляя пол, ограниченный обучением, а не транспортом. Это всё
энергетическое заявление, и оно честно: то же вычисление, пол на девять порядков
ниже, потому что нечего двигать и обратимое ядро бесплатно.

\begin{figure}[H]
\centering
\begin{tikzpicture}
\begin{axis}[
  width=12cm, height=6.4cm,
  xlabel={выполненные вычисления (операции, произв.\ единицы)},
  ylabel={энергия (log, Дж)},
  ymode=log, ymin=1e-21, ymax=1e-8,
  xmin=0, xmax=10, xtick=\empty,
  % легенда — в пустой средней полосе между верхней линией фон Неймана и нижними
  % полами Talos/Ландауэра: не перекрывает ни одну кривую; короткие
  % параллельные подписи, полное объяснение несёт подпись к рисунку.
  legend style={at={(0.5,0.52)}, anchor=center, font=\footnotesize,
                legend cell align=left, draw=gray!55, fill=white, rounded corners},
  grid=both, grid style={gray!15},
  every axis plot/.append style={very thick},
]
\addplot[plosblue,domain=0.2:10,samples=40] {2e-11*x + 3e-11};
\addlegendentry{пол фон Неймана ($\propto$ счёт)}
\addplot[talosamber,domain=0.2:10,samples=60] {1.4e-19*exp(-0.45*x) + 6e-21};
\addlegendentry{пол \textsc{Talos} (ограничен обучением)}
\addplot[plosgray,dashed,domain=0.2:10,samples=2] {2.87e-21};
\addlegendentry{пол Ландауэра $kT\ln2$}
\end{axis}
\end{tikzpicture}
\caption{Учёт энергии (log-шкала, схематично с машинно-проверенными якорями).
Пол фон Неймана растёт с объёмом вычислений, потому что каждая операция двигает
данные (синий). Пол \TALOS{} задан скоростью \emph{обучения}: высок при
сходимости, спадает к полу Ландауэра (серый пунктир), когда холон достигает
$\rhostar$ и перестаёт учиться (янтарный). Якоря:
$kT\ln2=2{,}87\times10^{-21}$~Дж; DRAM-слово $\approx2\times10^{-11}$~Дж.}
\label{fig:energy}
\end{figure}

\subsection{Метаболический пол и ось рабочей точки}
\label{subsec:metabolic}

Закон выше учитывает \emph{вычисление} (обратимое, свободное от Ландауэра) и
\emph{обучение} (спадающий член $\Regen$). Остаётся один член: цена просто
\emph{оставаться живым}. Диссипатор $\Diss$ непрерывно толкает $\Gam$ к
максимально смешанному состоянию тепловой смерти $I/7$; удержание состояния
против этой утечки не бесплатно.

\begin{claimbox}[Метаболический пол \status{Т}$+$\status{С} (T-273)]
В стационаре жизнеспособный холон ($P>2/7$, $\Gam\neq I/7$) экспортирует энтропию
со строго-положительной скоростью диссипатора $\dot S_{\Diss}>0$ (закон энтропии
T-271). По Ландауэру минимальная мощность \emph{поддержания}
\[
  P_{\mathrm{meta}} \;\geq\; f\cdot kT\ln 2\cdot \dot S_{\Diss}(\text{бит/тик}) \;>\; 0,
\]
\emph{активный, необратимый} двойник обратимого энергопола. Полная мощность
жизнеспособной когерентной машины поэтому есть $0$ (обратимый счёт) $+\
P_{\mathrm{meta}}$ (поддержание): она \textbf{не может работать бесплатно}. При
$300$~K пол одного холона --- фемто- до пиковатт (зависит от частоты тика), и он
измерим, ибо $\dot S_{\Diss}$ --- живая наблюдаемая телеметрии
(\S\ref{sec:platform}).
\end{claimbox}

Пол масштабируется с \textbf{удерживаемым порядком} (расстоянием $\ln 7 - S$ от
$I/7$): более когерентное, более сложное состояние сильнее толкается диссипатором
и стоит дороже удержать --- \emph{цена сложности}. Это обращает окно
жизнеспособности в явную \textbf{ось проектирования}, а не одну уставку.

\begin{claimbox}[Ось рабочей точки \status{С} (T-274)]
Жизнеспособное окно $(2/7,3/7]$ есть ось выразительность--эффективность, не одна
точка:
\begin{itemize}[leftmargin=1.4em,itemsep=1pt]
  \item \textbf{Экономный край $P\to 2/7$} --- эффективность-оптимальный:
        минимальная мощность поддержания, минимальный запас. Режим «выживания» ---
        наивысшая способность-на-ватт (для способности $C=\Phi\times R$), наименее
        робастный.
  \item \textbf{Богатый край $P\to 3/7$} (аттрактор T-124) ---
        способность-оптимальный: максимальное $C=2/3$ и максимальный динамический
        диапазон, при наивысшей мощности поддержания. Режим «процветания».
\end{itemize}
Динамика тянет машину к $P=3/7$ по умолчанию; инженер может гнать её экономнее к
$2/7$, меняя способность и запас на мощность. Оптимальная уставка --- \status{С}
(зависит от выбранной меры способности); \emph{робастна} же структура ---
монотонная цена по окну, два именованных режима на его краях. Машинно-проверено
(\texttt{t273\_metabolic.py}, \texttt{t274\_optimum.py}).
\end{claimbox}
